\documentclass[conference]{IEEEtran}

\usepackage[utf8]{inputenc}
\usepackage[T1]{fontenc}
\usepackage{mathptmx}
\usepackage{amsmath,amssymb}
\usepackage{array}
\usepackage{cite}
\usepackage{graphicx}
\usepackage[caption=false,font=footnotesize]{subfig}
\usepackage{url}
\usepackage{float}
\usepackage[section]{placeins}
\usepackage{xcolor}
\usepackage[colorlinks=true,linkcolor=green!45!black,citecolor=green!45!black,urlcolor=green!45!black]{hyperref}

\renewcommand{\topfraction}{0.95}
\renewcommand{\bottomfraction}{0.95}
\renewcommand{\textfraction}{0.05}
\renewcommand{\floatpagefraction}{0.85}
\renewcommand{\dbltopfraction}{0.95}
\renewcommand{\dblfloatpagefraction}{0.85}

\hypersetup{
pdfauthor={Nathnael Biresaw Mengist},
pdftitle={Residual Soft Actor-Critic with LQR, MPC, and H-Infinity Backbones for Robust DC Motor Speed Control},
pdfsubject={DC motor speed control with classical and reinforcement learning controllers},
pdfcreator={LaTeX},
pdfproducer={pdfTeX}
}


\graphicspath{
{../python_td3_experiments/outputs/experiment2_more_controllers_full_rl/}
{../python_td3_experiments/outputs/sac_balanced_10seeds_100k_final/}
{../python_td3_experiments/outputs/experiment4_condition_aware_10seeds_100k/}
{../python_td3_experiments/outputs/paper_direct_ddpg_50k/}
{../python_td3_experiments/outputs/residual_sac_curriculum_mpcwarm_10seeds_100k_best_comparison/}
{../python_td3_experiments/outputs/residual_sac_curriculum_mpcwarm_10seeds_100k_report_graphs/}
{../python_td3_experiments/outputs/residual_backbone_comparison_lqr_mpc_hinf/}
{./}
}

\begin{document}

\title{Residual Soft Actor-Critic with LQR, MPC, and H-Infinity Backbones for Robust DC Motor Speed Control}

\author{
\IEEEauthorblockN{Nathnael Biresaw Mengist}
\IEEEauthorblockA{
Department of Electrical and Computer Engineering\\
Kazakh-British Technical University\\
Almaty, Kazakhstan\\
Email: nathnaelbiresaw144@gmail.com}
}

\maketitle

\begin{abstract}
Classical controllers remain difficult baselines for reinforcement learning in motor-drive applications because LQR, MPC, and robust state-feedback controllers exploit accurate plant models and can use actuator limits aggressively. This paper presents a simulation benchmark for DC motor speed control comparing PID, LQR, LQI, MPC, H-infinity state feedback, direct deep reinforcement learning, and residual SAC controllers built on different classical backbones. Initial experiments show that directly trained SAC, TD3, and DDPG policies are not competitive with strong model-based baselines on tracking accuracy: direct SAC reduces control energy but leaves large steady-state error. The study therefore evaluates residual SAC controllers in which the learned policy supplies only a bounded voltage correction on top of LQR, MPC, or H-infinity state-feedback commands. All residual policies use the same MPC warm-start replay data, nominal-to-hard curriculum, reward design, ten random seeds, $100000$ training steps per seed, and validation-based best-checkpoint selection. Across twenty-four benchmark cases, standalone H-infinity obtains the lowest mean IAE, $15.6040$, while residual SAC-H-infinity obtains the lowest mean steady-state error, $5.9890 \pm 0.0524$. Residual SAC-LQR slightly improves LQR in mean IAE and SSE, and residual SAC-MPC improves standalone MPC, but neither dominates H-infinity on accumulated tracking error. The results show that reinforcement learning is most effective in this setting as a residual adaptive layer whose benefit depends strongly on the classical backbone.
\end{abstract}

\begin{IEEEkeywords}
DC motor control, reinforcement learning, soft actor-critic, residual reinforcement learning, LQR, MPC, H-infinity control, robust control.
\end{IEEEkeywords}

\section{Introduction}
DC motors are widely used in electric drives, robotics, conveyors, positioning systems, and laboratory control platforms. Their nominal dynamics are well understood and can be regulated effectively using classical methods such as proportional-integral-derivative (PID) control, linear quadratic regulation (LQR), linear quadratic integral (LQI) control, and model predictive control (MPC). For a linear time-invariant plant with a quadratic objective, LQR is not merely a convenient baseline; it is the analytic optimal feedback controller for the specified model and cost \cite{kalman1960,anderson1971}. This makes claims that reinforcement learning (RL) ``outperforms LQR'' difficult to interpret unless the comparison is carefully controlled.

Deep RL has nevertheless become attractive for motor control because real drives are not perfectly linear. Armature resistance changes with temperature, load inertia may vary, friction is nonlinear, magnetic saturation reduces torque production at high current, and voltage limits constrain achievable transients. These effects violate the assumptions under which a nominal LQR is optimal. The research question is therefore not whether model-free RL can rediscover LQR on a clean linear plant. The useful question is whether RL can improve robustness when the deployed plant differs from the nominal design model.

The first stage of this study tested direct RL controllers. TD3, DDPG, and SAC were trained and evaluated against PID, LQR, LQI, and MPC using the same DC motor model, the same voltage limit of $\pm 24$ V, the same reference signals, and the same nonlinear stress cases. That experiment showed an important negative result: direct SAC and TD3 did not beat LQR or MPC on tracking. SAC used much less energy, but this was caused by under-actuation and produced large final error. TD3 tracked better than direct SAC but remained weaker than LQR and MPC.

The next stage of the work therefore moved from pure RL to residual RL. The main idea is simple. Instead of letting the neural network control the motor completely by itself, a strong classical controller is kept in the loop and RL is asked to make only a small correction when needed. The first residual design used LQR as the backbone and SAC as the learning algorithm. After that result, the same training pipeline was extended to MPC and H-infinity state-feedback backbones. This addition is important because it tests whether the apparent benefit of residual SAC is a general RL advantage or whether it depends on the classical controller underneath the learned correction.

LQR was chosen as the first classical backbone for several specific reasons. First, the DC motor benchmark begins from a nominal state-space model, and LQR is the most natural optimal controller for that setting because it directly uses that model and provides an analytic full-state feedback law for a quadratic tracking-effort tradeoff. Second, LQR is much lighter and easier to embed inside a residual RL loop than MPC. MPC is also strong in this paper, but it solves an optimization problem online at every step, whereas LQR provides a fixed gain that is simple, fast, and stable to use as a baseline command inside the RL environment. Third, LQR is easier to interpret than PID or LQI as a residual backbone. PID does not provide the same full-state optimal structure, and LQI adds integral action that can partly hide the residual effect that the learned policy is supposed to correct. By contrast, LQR gives a clean nominal controller: strong enough to stabilize and track well, but simple enough that the extra contribution from SAC remains visible and meaningful. However, a strong journal paper should not stop at one backbone. MPC is included because it is a strong constraint-aware controller, and H-infinity state feedback is included because it is designed explicitly around disturbance attenuation and robustness. Comparing residual SAC-LQR, residual SAC-MPC, and residual SAC-H-infinity therefore makes the contribution more precise: the paper studies when residual RL helps, not just whether one hybrid controller can beat one classical baseline.

In practical terms, residual SAC with a classical backbone means the following. The backbone controller is responsible for the basic job of stabilizing the motor and following the reference using the nominal mathematical model. SAC does not replace that controller. Instead, SAC observes the motor state, the tracking error, and the operating condition, and then adds a limited extra voltage correction when the backbone command alone is not enough. In other words, the classical controller handles the part of the problem that is already well understood, and SAC learns only the leftover part of the problem: the residual error caused by saturation, disturbance, friction, or model mismatch.

This residual structure is important for two reasons. First, it is easier to train than a pure RL controller because the neural policy does not need to discover the entire control law from scratch. Second, it is more trustworthy than direct RL because the controller keeps the stabilizing behavior of LQR and uses learning only where classical nominal control is weakest. For that reason, the proposed method should be understood as a hybrid control strategy, not as a claim that RL alone is better than classical control in every case.

Mathematically, each residual controller uses
\begin{equation}
u = u_b + u_{\mathrm{SAC}},
\label{eq:residual_control}
\end{equation}
where $u_b$ is the nominal command from a backbone controller $b \in \{\mathrm{LQR}, \mathrm{MPC}, H_\infty\}$ and $u_{\mathrm{SAC}}$ is a bounded learned correction. This structure gives the controller a stabilizing model-based backbone and lets RL focus on nonlinear compensation, reference-dependent corrections, and robustness under stress.

The final residual SAC controller is not trained on one fixed task from start to finish. Instead, it is trained with a curriculum that gradually makes the problem harder. In the conducted experiment, the four curriculum stages have the following meanings. \emph{Nominal uniform} means training only on the nominal motor model while sampling the four reference scenarios uniformly. \emph{Robust uniform} means training on a broader distribution of plant conditions, including nominal, parameter changes, friction, saturation, and combined stress, while still sampling the reference scenarios uniformly. \emph{Hard uniform} shifts the plant-condition sampling more strongly toward the most difficult nonlinear cases, especially saturation and combined stress, while keeping scenario sampling uniform. Finally, \emph{Hard tracking} keeps the hard plant-condition distribution and also changes the reference-scenario sampling so that more training episodes emphasize the more difficult tracking cases, especially step responses with disturbance and sustained nonlinear tracking.

This curriculum design also explains an observation that can look counterintuitive at first: the training reward declines in later stages instead of increasing upward. The reason is not necessarily that the controller is becoming worse. The reason is that the task itself becomes much harder. In the saved training summary, the mean reward is approximately $372.63$ during nominal uniform training, then falls to about $-35.95$ in robust uniform, to about $-195.60$ in hard uniform, and to about $-260.65$ in hard tracking. These reward values are not directly comparable as if they came from the same fixed environment. Each later stage introduces more severe nonlinearities, more saturation, and harder reference-following demands, so the reward scale shifts downward even when the learned controller is improving in the sense that matters for the paper. For this reason, the paper evaluates final performance using external benchmark metrics such as IAE, SSE, control energy, and saturation fraction rather than relying on raw final reward alone.

The main contributions of this paper are:
\begin{itemize}
    \item A fair benchmark is built for PID, LQR, LQI, MPC, H-infinity state feedback, TD3, DDPG, direct SAC, and residual SAC variants. All controllers are tested on the same plant, voltage constraint, reference signals, and nonlinear stress conditions.
    \item A clear negative result is reported for direct RL. The paper shows, with quantitative results, that directly trained SAC and TD3 policies are not competitive with strong LQR and MPC baselines on the main tracking metrics, even when SAC achieves much lower control energy.
    \item A residual SAC framework is evaluated on three backbones: LQR, MPC, and H-infinity state feedback. The training pipeline is kept fixed across backbones using MPC warm-start data, curriculum learning, and validation-based best-checkpoint selection.
    \item A ten-seed backbone study shows that standalone H-infinity achieves the lowest mean IAE, while residual SAC-H-infinity achieves the lowest mean steady-state error. This gives a more honest and useful conclusion than claiming that RL universally outperforms classical control.
    \item The paper provides a practically useful conclusion: in this DC motor benchmark, reinforcement learning becomes most effective when it is used to complement a strong classical controller, and its value must be judged separately for IAE, final error, energy, and robustness.
\end{itemize}

The core finding of the paper is therefore not that reinforcement learning replaces classical control, but that residual SAC becomes useful when it is constrained to improve a strong backbone under nonlinear stress. What makes this study different from the previous papers is that it does not stop at a single hybrid controller or a single baseline. Instead, it compares direct RL, classical control, and residual SAC under one common benchmark and then extends the residual analysis across three backbones, LQR, MPC, and H-infinity, with the same data, curriculum, seeds, and checkpoint-selection rule. This makes it possible to distinguish a genuine backbone-dependent hybrid-control effect from a result that might otherwise appear to come only from one selected controller pair.

Taken together, the introduction frames the paper as a controlled test of where residual reinforcement learning adds value beyond already strong classical motor control.

\subsection{Related Work}
\subsubsection{Classical Motor Control}
PID control remains widely used in industrial motor drives because of its simplicity, interpretability, and ease of implementation \cite{astrom1995}. When an adequate model is available, however, state-space methods offer a more systematic route to performance shaping. LQR provides an analytically optimal state-feedback law for an infinite-horizon quadratic objective \cite{kalman1960,anderson1971}, while LQI adds integral action to reduce persistent tracking offset. MPC extends this model-based framework by optimizing control moves over a finite prediction horizon while explicitly enforcing constraints, albeit at greater computational cost \cite{mayne2000}. H-infinity control brings a complementary robust-control perspective by attenuating the worst-case influence of disturbances on regulated outputs \cite{doyle1989,zhou1996}. Together, these methods define a strong family of baselines for motor-drive applications.

Their relevance to the present study is straightforward. LQR and MPC are difficult learning benchmarks precisely because they exploit plant structure directly: LQR through an analytic nominal optimal law, and MPC through repeated constrained optimization. In electric-drive applications, where actuator limits, transient quality, and disturbance rejection are central design concerns, such controllers already solve a substantial part of the problem. For that reason, meaningful claims about RL-based improvement should be made relative to strong classical designs rather than to weakly tuned baselines alone.

At the same time, the classical-control literature makes clear that nominal design quality does not guarantee uniform robustness. Friction, dead zones, load changes, temperature-dependent resistance, and magnetic nonlinearities can all degrade performance when the implemented plant departs from the design model. Friction-compensation and disturbance-rejection studies have long documented this gap between nominal optimality and nonlinear operating reality \cite{armstrong1994}. More recent PMSM research has continued in this direction by combining model-based regulation with adaptive or disturbance-aware mechanisms, including RL-assisted ADRC variants and hybrid robust designs \cite{wang2023tec,wang2023tpe}. This literature is important here because it frames the proper comparison: the challenge for RL is not to outperform a weak conventional controller, but to add value beyond already capable model-based control.

\subsubsection{Deep Reinforcement Learning for Continuous Control}
Modern continuous-control RL emerged through successive developments in policy-gradient and actor-critic learning. Deterministic Policy Gradient established the continuous-action formulation \cite{silver2014}, and DDPG extended it with deep function approximation \cite{lillicrap2015}. TRPO and PPO introduced alternative policy-optimization strategies aimed at improved stability in high-dimensional settings \cite{schulman2015,schulman2017}. TD3 addressed important overestimation issues in deterministic actor-critic learning through clipped double critics, target policy smoothing, and delayed policy updates \cite{fujimoto2018}. SAC subsequently became one of the strongest general-purpose continuous-control methods by combining off-policy learning with entropy regularization \cite{haarnoja2018}. These developments, together with broader surveys and reviews, have made deep RL a standard reference point for robotics and nonlinear control problems \cite{kormushev2013,francoislavet2018,recht2019,motorreview2024}.

The appeal of these methods to motor control is obvious: they operate naturally in continuous action spaces, can represent nonlinear control laws, and do not require a fully explicit analytic controller. Yet the RL literature also offers a cautionary lesson. Strong performance in generic continuous-control benchmarks does not automatically translate into high-quality engineering control. Sample inefficiency, reward-design sensitivity, training instability, and mismatch between internal learning objectives and external performance criteria remain persistent concerns \cite{recht2019,francoislavet2018}. In motor-drive settings, where overshoot, steady-state error, actuator usage, and saturation all matter simultaneously, these issues can become even more pronounced. Accordingly, improvements in training reward should not be interpreted as sufficient evidence of superior closed-loop control quality.

\subsubsection{Residual Reinforcement Learning}
Residual RL combines a conventional controller with a learned correction rather than replacing the baseline altogether. The central idea is to let the model-based controller handle stabilization and nominal tracking while the learned policy focuses on the remaining mismatch, uncertainty, or nonlinear residual dynamics. This perspective is consistent with adaptive optimal-control views of RL \cite{lewis2012} and with results showing that RL can recover LQR-like behavior, albeit often with far greater sample requirements than direct analytic design \cite{fazel2018}. An important milestone in this line is the residual RL framework for robot control reported by Johannink \emph{et al.} \cite{johannink2019}. Related sim-to-real work has also emphasized the value of robust training procedures such as domain randomization when learned policies must operate under model mismatch \cite{tobin2017,peng2018}.

For control engineering, the attraction of residual RL is that it changes the role of learning from full replacement to targeted compensation. Rather than asking RL to solve stabilization, tracking, and robustness simultaneously, the method assumes that a reliable nominal controller already exists and asks the policy to improve the part of the problem that remains difficult. This is especially appealing in motor drives, where complete replacement of a classical controller may be hard to justify, but bounded adaptive corrections can still be valuable. Residual structures also improve interpretability: nominal behavior can be attributed to the baseline controller, while the learned action can be viewed as an additive compensation term. In practice, this often improves trainability and repeatability relative to pure model-free control.

This perspective is particularly relevant for the backbone choices studied in the present paper. MPC is already a strong constrained-control baseline because it uses an internal model and explicitly respects actuator limits \cite{mayne2000}. As a residual backbone, it offers a clear division of labor: the optimizer handles nominal predictive regulation, while the learned policy addresses the remaining mismatch caused by saturation, friction, or parameter variation. More broadly, work in robotics and motor control increasingly suggests that learning is most effective when it augments a structured controller rather than replacing it outright \cite{johannink2019,motorreview2024,review2025pmsm}. From this viewpoint, residual SAC-MPC is a natural extension of the hybrid-control literature, even if the exact combination remains comparatively underexplored.

An even stronger test arises when the backbone is robust H-infinity feedback. Because H-infinity control is explicitly designed around disturbance attenuation and worst-case robustness \cite{doyle1989,zhou1996}, it provides a more demanding baseline than nominal optimal control alone. This is important methodologically: claims about RL-based improvement become far more meaningful when the classical baseline already addresses uncertainty in a principled way. Recent motor-drive studies have increasingly moved toward such combinations of learning with disturbance-aware or robust control structures, including RL-assisted ADRC-style methods and hybrid robust designs for PMSM drives \cite{wang2023tec,wang2023tpe}. Although these studies are not identical to residual SAC-H-infinity, they support the same broader trend toward layered hybrid architectures in which robust classical feedback carries most of the stabilization burden and learning refines the remaining error.

Even so, an important gap remains. Prior residual RL work is conceptually strong, and prior motor-control RL work is increasingly rich empirically, but there are still relatively few studies that compare residual learning systematically across multiple high-quality backbones such as LQR, MPC, and H-infinity within one fixed benchmark. That gap matters because the usefulness of a learned residual may depend strongly on what the baseline already does well. A residual on top of nominal LQR may improve one portion of the error profile, whereas a residual on top of a robust H-infinity controller may chiefly affect final convergence rather than accumulated error. The present paper addresses this gap by holding the SAC training pipeline fixed and varying only the classical backbone, thereby sharpening the related-work question from ``does residual RL help?'' to ``which classical structure leaves the most useful room for a learned correction?'' \cite{johannink2019,mayne2000,doyle1989,zhou1996}.

\subsubsection{RL for Motor and Drive Control}
The application of RL to electric-drive control has expanded rapidly in recent years. For DC motors, prior work has explored Q-learning-based control, adaptive PI tuning, TD3-based PI adaptation, and real-time RL implementations on hardware \cite{rl_dcmotor_dqn,adaptive_pi_dc_motor,tufenkci2023,poudel2022,anugula2021,qin2018}. In PMSM drives, the literature is broader still, with reports on TD3-based vector control, speed optimization, deep RL current control, RL-assisted disturbance rejection, and RL-based servo regulation, alongside several recent review articles \cite{yin2023td3pmsm,hu2025td3pmsm,schindler2023iemdc,wang2023tec,wang2023tpe,zhao2023pmsm,wfhu2023multiscenario,dukkipati2026,review2025pmsm}. Collectively, these studies show that RL has become an important research direction in motor control, but they also reveal recurring methodological limitations: many comparisons rely on restricted baselines, focus on a narrow operating regime, or fail to distinguish clearly between tracking quality and actuator usage.

Within the smaller DC-motor literature, two patterns are especially visible. Some studies use RL to tune or adapt an existing PI or PID controller, thereby lowering the learning burden and improving practical stability \cite{adaptive_pi_dc_motor,tufenkci2023,qin2018}. Others adopt a more direct end-to-end RL formulation and report promising simulation or hardware performance \cite{rl_dcmotor_dqn,poudel2022,anugula2021}. Yet the scope of evaluation is often limited, making it difficult to determine whether the learned policy is broadly superior or merely advantageous under a particular reference signal, disturbance type, or baseline choice.

The PMSM literature reinforces the same point from a wider application base. A range of recent studies has shown that TD3-, DDPG-, and related DRL schemes can improve current regulation, speed optimization, disturbance rejection, and parameter adaptation in simulation and, in some cases, experimental settings \cite{yin2023td3pmsm,hu2025td3pmsm,schindler2023iemdc,zhao2023pmsm,wfhu2023multiscenario}. At the same time, review papers increasingly note that many of the most practically credible RL motor-drive solutions are not pure model-free replacements, but hybrid strategies that tune, optimize, or correct classical controllers \cite{review2025pmsm,motorreview2024}. This observation aligns closely with the present study.

Taken together, the literature suggests two requirements for a convincing RL motor-control benchmark. First, the learned controller should be compared against strong classical baselines such as LQR, MPC, and robust H-infinity feedback, not against PID alone. Second, the evaluation should separate distinct performance notions, especially accumulated tracking error, final convergence, and control effort. The present paper is structured around exactly these two requirements. It first documents the negative result that direct TD3, DDPG, and SAC do not outperform the strongest classical baselines in this benchmark, and then shows that residual SAC becomes more competitive precisely because it builds on rather than discards classical control structure.

In that sense, the related-work discussion motivates both the benchmark design and the hybrid residual direction adopted in the rest of the paper.

\section{Methodology}

\begin{figure}[!b]
\centering
\includegraphics[width=\columnwidth]{methodology_flow_diagram.pdf}
\caption{Methodological flow of the DC motor benchmark, controller design, residual SAC training, and evaluation pipeline.}
\label{fig:methodology_flow}
\end{figure}

\subsection{DC Motor Model}
\subsubsection{Nominal Dynamics}
The plant is an armature-controlled DC motor with state
\begin{equation}
x(t) = \begin{bmatrix} i_a(t) & \omega(t) \end{bmatrix}^{T},
\end{equation}
where $i_a$ is armature current and $\omega$ is angular speed. The input is armature voltage $u(t)$ and the output is speed. The nominal continuous-time model is
\begin{align}
\dot{i}_a &= -\frac{R_a}{L_a} i_a - \frac{K_e}{L_a}\omega + \frac{1}{L_a}u, \\
\dot{\omega} &= \frac{K_t}{J} i_a - \frac{B_m}{J}\omega - \frac{1}{J}\tau_L,
\end{align}
where $\tau_L$ is load torque. In state-space form,
\begin{align}
\dot{x} &= Ax + Bu + Ed, \\
y &= Cx,
\end{align}
with
\begin{equation}
A =
\begin{bmatrix}
-R_a/L_a & -K_e/L_a \\
K_t/J & -B_m/J
\end{bmatrix},
\quad
B =
\begin{bmatrix}
1/L_a \\
0
\end{bmatrix}.
\end{equation}

The nominal parameters used in simulation are shown in Table~\ref{tab:parameters}. Voltage is limited to $\pm 24$ V for every controller.

\begin{table}[!htbp]
\caption{Nominal DC Motor Parameters}
\label{tab:parameters}
\centering
\begin{tabular}{lcc}
\hline
Parameter & Symbol & Value \\
\hline
Armature resistance & $R_a$ & $2.59~\Omega$ \\
Armature inductance & $L_a$ & $0.58$ mH \\
Rotor inertia & $J$ & $5.69 \times 10^{-4}$ kg m$^2$ \\
Viscous friction & $B_m$ & $1.0 \times 10^{-6}$ N m s/rad \\
Torque constant & $K_t$ & $28.6 \times 10^{-3}$ \\
Back-EMF constant & $K_e$ & $28.6 \times 10^{-3}$ \\
Voltage limit & $u_{\max}$ & $24$ V \\
Episode duration & $T$ & $0.5$ s \\
Control step & $\Delta t_c$ & $10^{-3}$ s \\
\hline
\end{tabular}
\end{table}

\subsubsection{Perturbed and Nonlinear Conditions}
The benchmark evaluates all controllers on six plant conditions:
\begin{itemize}
    \item nominal plant;
    \item $R_a+50\%$ resistance increase;
    \item $J+50\%$ inertia increase;
    \item Coulomb friction with $F_c = 0.005$ N m;
    \item magnetic saturation with $\alpha = 0.15$;
    \item combined stress with resistance increase, friction, and saturation.
\end{itemize}

Magnetic saturation is modeled as a current-dependent torque constant,
\begin{equation}
K_{t,\mathrm{eff}}(i_a) = \frac{K_t}{1 + \alpha |i_a|}.
\end{equation}
Coulomb friction, a common nonlinear effect in mechanical drives \cite{armstrong1994}, is modeled as
\begin{equation}
\tau_f = F_c \operatorname{sign}(\omega),
\end{equation}
with zero friction torque at exactly zero speed.

\subsection{Controllers and Training Method}
\subsubsection{Classical Controllers}
The PID, LQR, LQI, MPC, and H-infinity controllers are designed on the nominal model and evaluated without retuning on perturbed cases. LQR uses the nominal state feedback law
\begin{equation}
u_{\mathrm{LQR}} = -Kx + N_r r,
\end{equation}
where $r$ is the speed reference and $N_r$ is a reference precompensator. The LQR design minimizes the infinite-horizon quadratic performance index
\begin{equation}
J_{\mathrm{LQR}} = \int_0^\infty (x^T Qx + u^T Ru)\,dt.
\end{equation}
In this expression, the matrix $Q \succeq 0$ weights the state-tracking objectives and $R \succ 0$ weights control effort. For the implemented controller,
\begin{equation}
Q = \mathrm{diag}(0.003,\ 2.8),
\qquad
R = [0.045].
\end{equation}
The LQR feedback gain is obtained from the continuous algebraic Riccati equation
\begin{equation}
A^T P + P A - P B R^{-1} B^T P + Q = 0,
\end{equation}
and
\begin{equation}
K = R^{-1} B^T P.
\end{equation}
The matrix $P$ is the positive semidefinite solution of the Riccati equation and defines the optimal cost-to-go for the nominal linear model. The term $-Kx$ provides the state-feedback action, while the precompensator $N_r$ improves reference tracking by scaling the command so that the closed-loop system reaches the desired steady-state output under the nominal model. In implementation,
\begin{equation}
N_r = -\big(C(A-BK)^{-1}B\big)^{-1}.
\end{equation}
This means that LQR handles two tasks simultaneously: it stabilizes the motor through feedback and it supplies a model-based feedforward term for reference tracking.

MPC uses the same nominal model structure and the same actuator limit, but it chooses the voltage by repeatedly solving a constrained prediction problem. The continuous motor model is first represented at the controller sampling instant as
\begin{align}
x_{k+1} &= A_d x_k + B_d u_k + E_d d_k, \\
y_k &= Cx_k,
\end{align}
where $A_d$, $B_d$, and $E_d$ are the discrete-time matrices, $d_k$ is the load disturbance, and $y_k$ is motor speed. At time $k$, MPC predicts future states over a horizon $N$ and minimizes
\begin{align}
J_{\mathrm{MPC}} =
&\sum_{j=0}^{N-1}
\big[
e_{k+j|k}^T Q_y e_{k+j|k}
+ x_{k+j|k}^T Q_x x_{k+j|k} \nonumber\\
&+ u_{k+j|k}^T R_u u_{k+j|k}
+ \Delta u_{k+j|k}^T R_{\Delta u}\Delta u_{k+j|k}
\big] \nonumber\\
&+ x_{k+N|k}^T Q_f x_{k+N|k},
\end{align}
where $e_{k+j|k}=y_{k+j|k}-r_{k+j}$ is the predicted tracking error.
The optimization is subject to
\begin{align}
x_{k+j+1|k} &= A_d x_{k+j|k}+B_d u_{k+j|k}, \\
-24 &\leq u_{k+j|k} \leq 24, \\
\Delta u_{k+j|k} &= u_{k+j|k}-u_{k+j-1|k}.
\end{align}
Here $Q_y$ penalizes speed-tracking error, $Q_x$ penalizes undesirable state magnitude, $R_u$ penalizes voltage usage, $R_{\Delta u}$ penalizes rapid voltage changes, and $Q_f$ is a terminal weight. Only the first optimized voltage is applied:
\begin{equation}
u_{\mathrm{MPC}}(k)=u_{k|k}^{\star}.
\end{equation}
At the next sample the horizon is shifted forward and the problem is solved again. This receding-horizon structure makes MPC a strong constraint-aware baseline because voltage saturation is handled directly inside the optimization, but it is computationally heavier than fixed-gain feedback.

The H-infinity baseline is implemented as robust state feedback with the load-torque channel treated as a disturbance input. In compact form, the disturbed plant is
\begin{equation}
\dot{x}=Ax+B_2u+B_1w,
\end{equation}
where $w$ represents external disturbance and modeling stress entering the mechanical dynamics. The regulated output is chosen as
\begin{equation}
z =
\begin{bmatrix}
Q_h^{1/2}x \\
R_h^{1/2}u
\end{bmatrix},
\end{equation}
so that the design penalizes both state error and control effort. The H-infinity objective is to make the worst-case energy gain from disturbance $w$ to regulated output $z$ smaller than a prescribed value $\gamma$:
\begin{equation}
\frac{\|z\|_2}{\|w\|_2} < \gamma,
\qquad
\|T_{wz}(s)\|_{\infty} < \gamma.
\end{equation}
This means that even if the disturbance takes an unfavorable shape, its effect on the weighted state and control response should remain bounded. A smaller $\gamma$ represents a more aggressive robustness requirement, while a larger $\gamma$ relaxes disturbance attenuation.

For full-state feedback $u=-K_{H_\infty}x+N_h r$, the gain is obtained from the continuous-time H-infinity Riccati condition
\begin{equation}
A^T P + PA + Q_h
-P\left(B_2R_h^{-1}B_2^T-\gamma^{-2}B_1B_1^T\right)P = 0,
\end{equation}
followed by
\begin{equation}
K_{H_\infty}=R_h^{-1}B_2^T P.
\end{equation}
The reference precompensator $N_h$ is computed in the same spirit as the LQR precompensator so that the nominal closed loop tracks a constant speed command. The implemented H-infinity controller uses the tuned weights
\begin{equation}
Q_h=\mathrm{diag}(0.001,\ 7.0),
\qquad
R_h=0.045,
\qquad
\gamma=50.
\end{equation}
This is not a full dynamic output-feedback H-infinity design; it is a robust full-state-feedback baseline intended to test whether disturbance-aware classical feedback is already strong enough to challenge residual RL. For a fair comparison, no controller is re-tuned separately for individual test cases.

\subsubsection{Direct RL Controllers}
Direct TD3, DDPG, and SAC policies output the full normalized motor-voltage command. The action is scaled to the voltage range and clipped:
\begin{equation}
u = \operatorname{clip}(u_{\max}a, -u_{\max}, u_{\max}), \quad a \in [-1,1].
\end{equation}
The observation vector contains normalized current, speed, reference speed, tracking error, integral error, error derivative, and load torque. Reward shaping penalizes absolute error, squared error, integral error, current, voltage magnitude, and voltage-rate variation, with small bonuses for tight tracking.

Because the final proposed method is based on SAC, its learning objective is described here in more detail. SAC learns a stochastic policy $\pi_\phi(a|s)$ and two soft Q-functions $Q_{\theta_1}(s,a)$ and $Q_{\theta_2}(s,a)$ from replay-buffer data. The goal is not only to maximize expected return, but also to maintain useful exploration through an entropy term. The soft objective can be written as
\begin{equation}
J(\pi) = \sum_{t=0}^{\infty} \mathbb{E}_{(s_t,a_t)\sim \rho_\pi}
\left[
r(s_t,a_t) + \alpha \mathcal{H}\big(\pi(\cdot|s_t)\big)
\right],
\end{equation}
where $\alpha$ is the entropy-temperature coefficient and $\mathcal{H}(\pi(\cdot|s))$ is the policy entropy. A larger entropy term encourages broader exploration, while a smaller entropy term makes the policy more deterministic.

The critic is trained using a soft Bellman target. For a transition $(s_t,a_t,r_t,s_{t+1})$, the target is
\begin{equation}
y_t =
r_t + \gamma
\left(
\min_{i\in\{1,2\}} Q_{\bar{\theta}_i}(s_{t+1},a_{t+1})
- \alpha \log \pi_\phi(a_{t+1}|s_{t+1})
\right),
\end{equation}
where $a_{t+1} \sim \pi_\phi(\cdot|s_{t+1})$ and $\bar{\theta}_i$ denotes target-network parameters. Each critic is then updated by minimizing
\begin{equation}
J_Q(\theta_i) =
\mathbb{E}\left[
\frac{1}{2}\big(Q_{\theta_i}(s_t,a_t)-y_t\big)^2
\right].
\end{equation}

The actor is updated to maximize the expected soft value, or equivalently to minimize
\begin{equation}
J_\pi(\phi) =
\mathbb{E}_{s_t \sim \mathcal{D},\, a_t \sim \pi_\phi}
\left[
\alpha \log \pi_\phi(a_t|s_t) - Q_{\theta_1}(s_t,a_t)
\right].
\end{equation}
This actor objective has a simple interpretation. The policy is encouraged to choose actions with high long-term value according to the critic, while also preserving enough entropy to avoid premature collapse to poor deterministic actions.

In this work, the SAC reward is not a single tracking-error term. It is a shaped control reward designed to reflect motor-control quality. The implemented reward includes: tracking bonuses near the reference; penalties on absolute and squared tracking error; an integral-error penalty to discourage accumulated offset; penalties on voltage magnitude, current magnitude, and voltage-rate change; and additional penalties for excessive saturation or ineffective high-voltage use. At the final time step, the reward adds extra emphasis on small final error. This reward design is important because naive low-energy rewards can produce controllers that save voltage by failing to track the reference.

\subsubsection{Residual SAC with Classical Backbones}
The residual controller keeps a classical baseline in the loop and lets SAC learn only a correction:
\begin{equation}
u = \operatorname{clip}(u_b + \Delta u_{\mathrm{SAC}}, -24, 24),
\end{equation}
where $u_b$ is supplied by LQR, MPC, or H-infinity state feedback.
The residual action is bounded to $\pm 8$ V. If the SAC policy outputs a normalized residual action $a_{\mathrm{SAC}} \in [-1,1]$, then the learned correction is
\begin{equation}
\Delta u_{\mathrm{SAC}} = 8\, a_{\mathrm{SAC}}.
\end{equation}
The final applied voltage is therefore
\begin{equation}
u = \operatorname{clip}(u_b + 8 a_{\mathrm{SAC}}, -24, 24).
\end{equation}
This bound prevents the policy from destroying the stabilizing behavior of the backbone while still allowing meaningful correction under nonlinear stress. Conceptually, the residual policy is trained to answer a narrower question than direct RL: ``Given that a strong nominal command is already available, what small extra correction should be added to reduce the remaining tracking error under nonlinear or mismatched conditions?''

The residual SAC design changes the learning problem in an important way. Direct SAC must discover both stabilization and fine tracking. Residual SAC does not start from zero; it starts from a meaningful classical baseline. This usually improves sample efficiency and reduces the chance that the policy learns an energy-conservative but inaccurate control strategy. In this sense, the residual formulation acts as a structured inductive bias that guides RL toward useful corrections instead of forcing it to rediscover the entire controller. The same SAC architecture, reward, curriculum, warm-start replay data, and checkpoint-selection rule are used for residual SAC-LQR, residual SAC-MPC, and residual SAC-H-infinity. Therefore, differences among the three residual controllers can be interpreted mainly as the effect of the selected classical backbone.

Three training improvements are used:
\begin{itemize}
    \item \textbf{MPC warm-start:} an expert replay buffer is initialized with MPC-generated transitions.
    \item \textbf{Curriculum learning:} training progresses from nominal cases to robust and hard nonlinear tracking cases.
    \item \textbf{Best-checkpoint selection:} saved checkpoints are evaluated on validation cases and the best checkpoint per seed is selected by validation performance.
\end{itemize}

Each final residual SAC experiment uses ten random seeds, $100000$ training steps per seed, a chunk size of $50000$ steps, and best-checkpoint selection.

\subsection{Experimental Protocol}
\subsubsection{Reference Signals}
Each controller is evaluated on four reference scenarios:
\begin{itemize}
    \item step command;
    \item step command with load disturbance;
    \item ramp command;
    \item sinusoidal command.
\end{itemize}
Combining four references with six plant conditions gives twenty-four test cases.

\subsubsection{Performance Metrics}
The experiments use several metrics because no single number fully describes a motor controller. Let
\begin{equation}
e(t)=r(t)-\omega(t)
\end{equation}
be the speed-tracking error. The main continuous-time metrics are
\begin{align}
\mathrm{IAE} &= \int_0^T |e(t)|\,dt, \\
\mathrm{ISE} &= \int_0^T e^2(t)\,dt, \\
\mathrm{SSE} &= |e(T)|, \\
E_u &= \int_0^T u^2(t)\,dt,
\end{align}
where $T$ is the test duration. In the numerical simulation these integrals are computed using the sampled approximation
\begin{align}
\mathrm{IAE} &\approx \Delta t \sum_{k=0}^{N_T} |e_k|, \\
\mathrm{ISE} &\approx \Delta t \sum_{k=0}^{N_T} e_k^2, \\
E_u &\approx \Delta t \sum_{k=0}^{N_T} u_k^2.
\end{align}
IAE is the primary tracking metric in this paper because it measures total accumulated tracking error over the whole response. A controller with lower IAE keeps the motor closer to the reference for more of the experiment. ISE is also reported because it squares the error and therefore penalizes large transient deviations more strongly than IAE. SSE measures the final remaining error; it is important for speed regulation because a controller can have acceptable transient behavior but still finish with a nonzero offset.

Control energy $E_u$ measures how aggressively the controller uses voltage. A lower value means less squared voltage effort, but it must be interpreted together with IAE and SSE. For example, the direct SAC controller has low energy mainly because it under-actuates the motor and leaves large tracking error. Therefore, low energy alone is not considered a success.

The saturation fraction is
\begin{equation}
S_{\mathrm{sat}} =
\frac{1}{N_T+1}
\sum_{k=0}^{N_T}
\mathbf{1}\left(|u_k| \geq u_{\max}-\epsilon\right),
\end{equation}
where $\mathbf{1}(\cdot)$ is the indicator function and $\epsilon$ is a small numerical tolerance. This metric is important because a controller that spends most of the test near $\pm 24$ V may be fragile on real hardware, even if its tracking error is low.

For step-response interpretation, the study also records time-domain diagnostics. Rise time $t_r$ is the time required for the response to move from $10\%$ to $90\%$ of the final reference value. Settling time $t_s$ is the first time after which the speed remains inside a small tolerance band around the reference, typically $2\%$ or $5\%$. Percent overshoot is
\begin{equation}
M_p =
\frac{\omega_{\max}-r_{\mathrm{final}}}{|r_{\mathrm{final}}|}
\times 100\%,
\end{equation}
when the response exceeds the final reference. These time-domain metrics are secondary to IAE and SSE in the tables, but they help explain whether a controller is fast, oscillatory, slow, or overly conservative.

Overall controller quality is therefore judged using a hierarchy. IAE is used as the main accumulated tracking score, SSE is used to assess final convergence, $E_u$ and saturation fraction are used to check actuator effort, and rise time, settling time, and overshoot are used to interpret transient shape.

\textit{Taken together, this methodology defines a common model, controller set, training pipeline, and evaluation protocol so that all reported comparisons remain consistent and reproducible.}

\section{Results and Discussion}
\subsection{Experiment 1: Direct Controller Comparison}
The first comparison evaluates direct classical and direct RL controllers. Fig.~\ref{fig:controller_heatmap} shows the controller comparison heatmap, while Fig.~\ref{fig:direct_iae_bar} shows mean IAE and Fig.~\ref{fig:direct_energy_bar} shows mean control energy. These figures are useful because they reveal not only which controller is best on average, but also whether a low-energy controller is genuinely efficient or simply failing to track the reference well.

\begin{figure}[!htbp]
\centering
\includegraphics[width=\columnwidth]{controller_comparison_heatmap.png}
\caption{Controller comparison heatmap for the direct-controller benchmark.\protect\\ Lower tracking-error metrics are better.}
\label{fig:controller_heatmap}
\end{figure}

\begin{figure}[!htbp]
\centering
\includegraphics[width=\columnwidth]{mean_iae_by_controller.png}
\caption{Direct-controller benchmark: mean IAE by controller.\protect\\ LQR and MPC obtain the lowest mean IAE.}
\label{fig:direct_iae_bar}
\end{figure}

\begin{figure}[!htbp]
\centering
\includegraphics[width=\columnwidth]{mean_control_energy_by_controller.png}
\caption{Direct-controller benchmark: mean control energy by controller.\protect\\ Direct SAC obtains low energy by under-tracking the reference.}
\label{fig:direct_energy_bar}
\end{figure}

\begin{table}[!htbp]
\caption{Direct Controller Benchmark}
\label{tab:direct_results}
\centering
\begin{tabular}{lccc}
\hline
Controller & Mean IAE & Energy & SSE \\
\hline
LQR & 15.71 & 202.10 & 6.73 \\
MPC & 15.74 & 201.78 & 6.81 \\
PID & 19.81 & 147.28 & 16.93 \\
TD3 & 20.32 & 148.26 & 20.02 \\
LQI & 21.91 & 141.68 & 23.22 \\
DDPG & 26.24 & 71.24 & 35.62 \\
SAC & 32.78 & 23.95 & 58.32 \\
\hline
\end{tabular}
\end{table}

Table~\ref{tab:direct_results} reports the mean direct-controller results. LQR and MPC are effectively tied for the best tracking performance, with mean IAE values of $15.71$ and $15.74$, respectively, and mean SSE values of $6.73$ and $6.81$. Numerically, this indicates that both classical optimal controllers keep total tracking error low across the full benchmark and finish close to the desired speed at the end of the test horizon. Fig.~\ref{fig:controller_heatmap} supports this reading visually by showing that low-error performance appears across most condition-scenario combinations, not only in one special case.

The direct RL results reveal a clear ranking. TD3 is the strongest of the pure RL controllers, DDPG is weaker, and SAC is the weakest on tracking although it uses the least energy. Direct TD3 reaches a mean IAE of $20.32$ and a mean SSE of $20.02$. Relative to LQR, this corresponds to roughly $29.3\%$ larger accumulated tracking error and about $197\%$ larger final error. Direct TD3 therefore learns a workable controller, but not one that matches the strongest classical baselines. Direct DDPG is weaker than TD3, with mean IAE $26.24$ and mean SSE $35.62$. Relative to TD3, DDPG has about $29.1\%$ larger IAE and about $77.9\%$ larger final error, even though its mean control energy drops from $148.26$ to $71.24$. This shows that DDPG also trades away too much tracking quality in exchange for reduced actuator use. The direct SAC result is even more revealing. Although direct SAC reduces mean control energy to only $23.95$, compared with $202.10$ for LQR, its mean IAE increases to $32.78$ and its mean SSE increases to $58.32$. Relative to LQR, that is about $108.6\%$ larger IAE and about $766\%$ larger final error. This means the low-energy behavior of direct SAC is not an efficiency improvement in the control-theoretic sense; rather, it reflects persistent under-actuation and poor reference tracking.

The direct SAC and TD3 learning curves in Fig.~\ref{fig:direct_sac_training_curve} and Fig.~\ref{fig:direct_td3_training_curve} must also be interpreted carefully. Both policies improve during training, so the negative result is not that learning failed completely. The real conclusion is that reward improvement during training was not enough to produce a controller that was competitive on the evaluation metrics used in this paper. In other words, the RL agents learned something, but what they learned was still inferior to the tracking behavior delivered by strong nominal LQR and MPC designs. The DDPG learning curve in Fig.~\ref{fig:ddpg_training_curve} shows the same overall pattern. Its episode reward improves from poor initial behavior during nominal direct training, but the final benchmark performance remains weak. This matters because it shows that training-progress improvement alone is not a sufficient sign of controller quality. A policy can learn steadily in reward terms and still remain significantly inferior to the best classical baselines when judged by IAE and steady-state error.

This direct-controller study therefore establishes the paper's first main result: pure RL is not yet competitive with the strongest classical baselines in this benchmark.

\begin{figure}[!htbp]
\centering
\includegraphics[width=\columnwidth]{sac_learning_curves.png}
\caption{Direct SAC learning curve in the initial benchmark.\protect\\ Reward improves during training, but final tracking remains uncompetitive against LQR and MPC.}
\label{fig:direct_sac_training_curve}
\end{figure}

\begin{figure}[!htbp]
\centering
\includegraphics[width=\columnwidth]{td3_learning_curves.png}
\caption{Direct TD3 learning curve in the initial benchmark.\protect\\ TD3 learns stronger tracking than SAC, but still finishes with larger IAE and SSE than the classical baselines.}
\label{fig:direct_td3_training_curve}
\end{figure}
\begin{figure}[!htbp]
\centering
\includegraphics[width=\columnwidth]{ddpg_learning_curves.png}
\caption{Direct DDPG learning curve for nominal 50k-step training.\protect\\ The reward trend improves from poor initial behavior, but the final direct DDPG controller still remains weak in benchmark tracking performance, with mean IAE $26.24$ and mean SSE $35.62$, which are both much worse than direct TD3 and much worse than LQR/MPC.}
\label{fig:ddpg_training_curve}
\end{figure}

\begin{table*}[!t]
\caption{Direct-controller winner table by condition and scenario}
\label{tab:direct_winner_table}
\centering
\scriptsize
\setlength{\tabcolsep}{4pt}
\renewcommand{\arraystretch}{1.08}
\begin{tabular}{lllll}
\hline
Condition & Scenario & Best IAE Controller & Best Energy Controller & Best SSE Controller \\
\hline
combined\_stress & ramp & LQR & SAC & LQR \\
combined\_stress & sine & LQR & SAC & LQR \\
combined\_stress & step\_load\_disturbance & LQR & SAC & LQR \\
combined\_stress & step\_nominal & LQR & SAC & LQR \\
friction & ramp & LQR & LQI & LQR \\
friction & sine & LQR & SAC & LQR \\
friction & step\_load\_disturbance & LQR & SAC & LQR \\
friction & step\_nominal & LQR & SAC & LQR \\
j\_plus\_50 & ramp & LQR & SAC & LQR \\
j\_plus\_50 & sine & LQR & SAC & TD3\_loaded \\
j\_plus\_50 & step\_load\_disturbance & LQR & SAC & LQR \\
j\_plus\_50 & step\_nominal & LQR & SAC & LQR \\
nominal & ramp & DDPG & LQI & DDPG \\
nominal & sine & LQR & SAC & LQR \\
nominal & step\_load\_disturbance & LQR & SAC & LQR \\
nominal & step\_nominal & LQR & SAC & LQR \\
ra\_plus\_50 & ramp & LQR & SAC & LQR \\
ra\_plus\_50 & sine & LQR & SAC & LQR \\
ra\_plus\_50 & step\_load\_disturbance & LQR & SAC & LQR \\
ra\_plus\_50 & step\_nominal & LQR & SAC & LQR \\
saturation & ramp & LQR & SAC & LQR \\
saturation & sine & LQR & SAC & LQR \\
saturation & step\_load\_disturbance & LQR & SAC & LQR \\
saturation & step\_nominal & LQR & SAC & LQR \\
\hline
\end{tabular}
\end{table*}
\FloatBarrier

\subsection{Experiment 2: Residual SAC-LQR}
This experiment evaluates the residual SAC-LQR controller against LQR, MPC, PID, LQI, and TD3. Table~\ref{tab:final_results} gives the overall mean results across the twenty-four benchmark cases for this residual SAC-LQR study. Residual SAC-LQR obtains the best mean IAE and the best mean SSE. The IAE improvement over LQR is small, $0.075\%$, but it is consistent across the ten selected checkpoints, with a seed standard deviation of only $0.0170$. The SSE improvement is larger: $7.52\%$ relative to LQR and $8.51\%$ relative to MPC.

The numerical interpretation of Table~\ref{tab:final_results} is important. Residual SAC-LQR does not dominate every metric. Its mean ISE, $1089.52 \pm 0.18$, is slightly worse than the $1086.05$ of LQR and the $1086.23$ of MPC. This means the hybrid controller is not universally best in every transient measure. However, it achieves the lowest overall IAE and the lowest final error while keeping energy and saturation almost unchanged relative to the strongest classical baselines. The practical meaning is that the learned residual improves useful tracking quality without asking the motor to work much harder.

Fig.~\ref{fig:overall_iae_results} and Fig.~\ref{fig:overall_sse_results} add an important robustness interpretation. The residual SAC-LQR error bars are very small, especially when compared with direct TD3. For IAE, residual SAC-LQR has a seed standard deviation of only $0.0170$, whereas TD3 has $0.7287$. For SSE, residual SAC-LQR has a seed standard deviation of $0.0992$, whereas TD3 has $3.8801$. This indicates that the residual design is not only more accurate on average but also far more repeatable across random seeds, which strengthens the reliability of the reported improvement.

This makes the residual SAC-LQR result more than a marginal numerical gain; it is also a repeatable hybrid-control improvement.

\begin{table}[!tbp]
\caption{Ten-seed residual SAC-LQR benchmark using best selected checkpoints}
\label{tab:final_results}
\centering
\scriptsize
\setlength{\tabcolsep}{3pt}
\renewcommand{\arraystretch}{1.1}
\begin{tabular}{lcccccc}
\hline
Metric & Residual SAC-LQR & LQR & MPC & PID & TD3 & LQI \\
\hline
Family & RL hybrid & Classical & Classical & Classical & RL & Classical \\
Seeds & 10 & -- & -- & -- & 10 & -- \\
IAE & $\mathbf{15.7019 \pm 0.0170}$ & 15.7137 & 15.7381 & 19.8084 & $20.4985 \pm 0.7287$ & 21.9115 \\
ISE & $1089.52 \pm 0.18$ & 1086.05 & 1086.23 & 1217.58 & $1382.81 \pm 55.31$ & 1356.54 \\
Energy & $202.45 \pm 0.22$ & 202.10 & 201.78 & 147.28 & $154.15 \pm 12.96$ & 141.68 \\
SSE & $\mathbf{6.2263 \pm 0.0992}$ & 6.7328 & 6.8053 & 16.9282 & $20.9578 \pm 3.8801$ & 23.2201 \\
Saturation & $0.6131 \pm 0.0016$ & 0.6147 & 0.6132 & 0.2768 & $0.1088 \pm 0.0661$ & 0.3493 \\
\hline
\end{tabular}
\end{table}
\begin{figure}[!htbp]
\centering
\includegraphics[width=\columnwidth]{overall_mean_iae_with_seed_std.png}
\caption{Overall mean IAE for the residual SAC-LQR benchmark.\protect\\ Error bars are seed standard deviations for learned controllers.}
\label{fig:overall_iae_results}
\end{figure}

\begin{figure}[!htbp]
\centering
\includegraphics[width=\columnwidth]{overall_mean_sse_with_seed_std.png}
\caption{Overall mean SSE for the residual SAC-LQR benchmark.\protect\\ Error bars are seed standard deviations for learned controllers.}
\label{fig:overall_sse_results}
\end{figure}
\FloatBarrier
\subsection{Hard Nonlinear Cases}
The hardest subset consists of saturation and combined-stress conditions. Table~\ref{tab:hard_results} shows that LQR has a slightly lower hard-case mean IAE than residual SAC-LQR, but residual SAC-LQR has lower SSE and lower control energy. This nuance is important: the residual RL controller is not uniformly superior in every metric. Its strongest advantage is final accuracy and ramp tracking, while LQR remains highly competitive on total tracking error in the most saturated cases.

The hard-case IAE difference is numerically very small. LQR obtains $22.7195$, while residual SAC-LQR obtains $22.7310 \pm 0.0319$. The absolute gap is only about $0.0115$, which is roughly $0.05\%$ of the LQR value. By contrast, the hard-case SSE difference is more meaningful: residual SAC-LQR reduces mean final error from $18.2369$ to $16.9766$, which is about a $6.9\%$ improvement. It also lowers hard-case mean control energy slightly, from $284.31$ to $283.49$. Therefore, Fig.~\ref{fig:hard_cases_iae} and Fig.~\ref{fig:hard_cases_energy} should be read as showing near parity on accumulated hard-case error, with a moderate advantage for the residual controller in final convergence quality.

This hard-case view therefore clarifies that the residual gain is selective and is strongest in final convergence rather than in universal dominance across all metrics.

\begin{table}[!htbp]
\caption{Hard-Case Mean Results}
\label{tab:hard_results}
\centering
\resizebox{\columnwidth}{!}{%
\begin{tabular}{lccc}
\hline
Controller & IAE & Energy & SSE \\
\hline
LQR & $\mathbf{22.7195}$ & 284.31 & 18.2369 \\
Residual SAC-LQR & $22.7310 \pm 0.0319$ & $\mathbf{283.49 \pm 0.53}$ & $\mathbf{16.9766 \pm 0.0645}$ \\
MPC & 22.7398 & 283.97 & 18.2789 \\
LQI & 25.5724 & 238.08 & 26.0779 \\
PID & 25.5877 & 207.28 & 28.3537 \\
TD3 & $27.0923 \pm 0.9596$ & $146.40 \pm 16.86$ & $33.4426 \pm 3.6213$ \\
\hline
\end{tabular}
}
\end{table}

\begin{figure}[!htbp]
\centering
\includegraphics[width=\columnwidth]{hard_cases_mean_iae_with_seed_std.png}
\caption{Hard nonlinear case comparison: mean IAE.\protect\\ LQR is slightly better on hard-case IAE.}
\label{fig:hard_cases_iae}
\end{figure}

\begin{figure}[!htbp]
\centering
\includegraphics[width=\columnwidth]{hard_cases_mean_energy_with_seed_std.png}
\caption{Hard nonlinear case comparison: mean control energy.\protect\\ Residual SAC-LQR uses slightly less hard-case energy.}
\label{fig:hard_cases_energy}
\end{figure}
\FloatBarrier
\subsection{Winner Analysis}
Across the twenty-four benchmark cases, LQR wins IAE in eighteen cases and residual SAC-LQR wins IAE in six cases. The six residual SAC-LQR wins occur in all ramp-tracking cases: nominal, resistance change, inertia change, friction, saturation, and combined stress. LQR dominates the step and sine cases.

\begin{figure}[!htbp]
\centering
\includegraphics[width=\columnwidth]{winner_count_bars.png}
\caption{Winner counts for the residual SAC-LQR benchmark.\protect\\ Residual SAC-LQR wins all ramp IAE cases, while LQR wins most step and sinusoidal cases.}
\label{fig:winner_counts}
\end{figure}

The winner pattern indicates that the residual policy mainly improves reference-following behavior where feedforward mismatch matters. For step and sinusoidal references, the nominal LQR remains extremely strong because high-gain feedback and the reference precompensator already provide accurate tracking. The winner table makes this point very concrete: the residual controller does not win everywhere, but it wins systematically in the entire ramp subset. This suggests that the learned residual is especially useful when the reference is continuously changing and the nominal feedforward structure no longer captures the true nonlinear plant behavior well enough.

The case-by-case analysis therefore sharpens the interpretation of where the residual policy is actually useful.
\FloatBarrier

\subsection{Training Curve Interpretation}
Fig.~\ref{fig:residual_sac_learning_curves} first shows the residual SAC learning curves across the selected training seeds. The individual curves vary because each seed experiences different replay-buffer trajectories and stochastic policy updates, but the overall pattern is consistent: reward improves early, then spreads as the curriculum introduces harder nonlinear operating cases. This is an important result by itself. It shows that the residual SAC controller is trainable and repeatable across seeds, but it also shows that learning-curve height alone is not the correct criterion for final controller quality. The final controller is selected using validation metrics, not simply the largest final reward.

Fig.~\ref{fig:curriculum_training} then shows the reward trajectory together with curriculum stages. The final reward becomes negative because the task distribution changes from nominal operation to harder nonlinear tracking. Therefore, final episode reward is not used as the main controller-quality measure. Validation IAE, SSE, control energy, and saturation fraction are more reliable for comparing final controllers.

This graph must be interpreted stage by stage rather than as one stationary reward history. Early in training, easier nominal or mildly perturbed cases are more common, so higher reward is easier to obtain. Later, the curriculum introduces saturation and combined-stress tracking cases, which are intrinsically harder and also penalize ineffective high-voltage behavior more strongly. As a result, the absolute reward becomes more negative even while the final selected controller becomes better on the benchmark metrics. This is why the paper relies on external evaluation metrics rather than raw final reward when judging controller quality.

This subsection therefore supports the choice of validation metrics and checkpoint selection as the main basis for final comparison.

\begin{figure}[!htbp]
\centering
\includegraphics[width=\columnwidth]{training_learning_curves_all_seeds.png}
\caption{Residual SAC-LQR learning curves across the selected training seeds.\protect\\ The curves show consistent learning progress, but the spread across seeds confirms that training reward alone is not sufficient for final model selection. This is why the study uses best-checkpoint validation on benchmark metrics rather than the last training checkpoint only.}
\label{fig:residual_sac_learning_curves}
\end{figure}

\begin{figure}[!htbp]
\centering
\includegraphics[width=\columnwidth]{training_reward_with_curriculum_stages.png}
\caption{Training reward with curriculum stages.\protect\\ Later stages introduce harder nonlinear tracking cases, so reward values are not directly comparable across the entire training run.}
\label{fig:curriculum_training}
\end{figure}
\FloatBarrier

\subsection{Experiment 3: Residual SAC-MPC and Residual SAC-H-Infinity}
After the residual SAC-LQR experiment, the same residual training method was repeated with MPC and H-infinity backbones. This experiment is important because it tests whether the learned residual is genuinely improving robust control behavior or only correcting weaknesses of one selected LQR baseline. The training recipe was not changed: each residual controller used the same reward, observation structure, curriculum stages, MPC warm-start replay data, ten random seeds, $100000$ steps per seed, and validation-based best-checkpoint selection. Only the backbone command $u_b$ was changed.

Table~\ref{tab:backbone_results} gives the main numerical result. The strongest conclusion is deliberately nuanced. Standalone H-infinity obtains the lowest mean IAE, $15.6040$, which means it has the best accumulated tracking accuracy across the full benchmark. Residual SAC-H-infinity obtains the lowest mean SSE, $5.9890 \pm 0.0524$, which means the learned correction improves final convergence accuracy even though it does not reduce total IAE below the robust H-infinity baseline. Residual SAC-LQR remains better than standalone LQR in both mean IAE and SSE. Residual SAC-MPC also improves standalone MPC, reducing mean IAE from $15.7381$ to $15.7168$ and reducing SSE from $6.8053$ to $6.3720$. However, residual SAC-MPC is not stronger than residual SAC-LQR or residual SAC-H-infinity. This shows that the backbone choice matters.

\begin{table*}[!tbp]
\caption{Backbone Comparison Across Twenty-Four Benchmark Cases}
\label{tab:backbone_results}
\centering
\scriptsize
\begin{tabular}{lccclp{0.25\textwidth}}
\hline
Controller & Seeds & IAE & SSE & Energy & Interpretation \\
\hline
H-infinity & -- & $\mathbf{15.6040}$ & 6.4062 & 203.54 & Best accumulated tracking error \\
Residual SAC-H-infinity & 10 & $15.6279 \pm 0.0318$ & $\mathbf{5.9890 \pm 0.0524}$ & $203.83 \pm 0.26$ & Best final error \\
Residual SAC-LQR & 10 & $15.7019 \pm 0.0170$ & $6.2263 \pm 0.0992$ & $202.45 \pm 0.22$ & Improves LQR and MPC \\
LQR & -- & 15.7137 & 6.7328 & 202.10 & Strong nominal baseline \\
Residual SAC-MPC & 10 & $15.7168 \pm 0.0598$ & $6.3720 \pm 0.3060$ & $202.13 \pm 0.34$ & Improves MPC but not H-infinity \\
MPC & -- & 15.7381 & 6.8053 & 201.78 & Strong constraint-aware baseline \\
\hline
\end{tabular}
\end{table*}
Fig.~\ref{fig:backbone_overall_iae} and Fig.~\ref{fig:backbone_overall_sse} visualize the same result. On the IAE plot, the H-infinity bar is lowest, so the paper should not claim that residual SAC defeats every classical controller. On the SSE plot, residual SAC-H-infinity is lowest, showing where the learned residual is most useful. Numerically, residual SAC-H-infinity reduces SSE by about $6.51\%$ relative to standalone H-infinity, from $6.4062$ to $5.9890$, while increasing energy only slightly. This is a meaningful result because steady-state error is directly visible in speed regulation applications.

The hard nonlinear subset in Table~\ref{tab:backbone_hard_results} and Fig.~\ref{fig:backbone_hard_iae} and Fig.~\ref{fig:backbone_hard_sse} confirms the same pattern under saturation and combined stress. H-infinity remains best on hard-case IAE, with $22.6280$. Residual SAC-H-infinity is very close in hard-case IAE, $22.6587 \pm 0.0390$, but has the best hard-case SSE, $16.8241 \pm 0.0778$. This means the residual SAC correction helps the final tracking point in the hardest nonlinear cases, but it does not eliminate the accumulated transient-error advantage of the robust H-infinity design.

\begin{table}[!htbp]
\caption{Backbone Comparison on Hard Nonlinear Cases}
\label{tab:backbone_hard_results}
\centering
\resizebox{\columnwidth}{!}{%
\begin{tabular}{lccc}
\hline
Controller & IAE & SSE & Energy \\
\hline
H-infinity & $\mathbf{22.6280}$ & 18.0499 & 285.84 \\
Residual SAC-H-infinity & $22.6587 \pm 0.0390$ & $\mathbf{16.8241 \pm 0.0778}$ & $285.49 \pm 0.41$ \\
LQR & 22.7195 & 18.2369 & 284.31 \\
Residual SAC-LQR & $22.7310 \pm 0.0319$ & $16.9766 \pm 0.0645$ & $\mathbf{283.49 \pm 0.53}$ \\
MPC & 22.7398 & 18.2789 & 283.97 \\
Residual SAC-MPC & $22.7409 \pm 0.0533$ & $16.9963 \pm 0.1079$ & $283.72 \pm 0.75$ \\
\hline
\end{tabular}
}
\end{table}

\begin{figure}[!b]
\centering
\subfloat[Overall backbone comparison: mean IAE. H-infinity gives the lowest accumulated tracking error.\label{fig:backbone_overall_iae}]{
\includegraphics[width=0.88\columnwidth]{overall_backbone_iae.png}}\\[-1mm]
\subfloat[Overall backbone comparison: mean SSE. Residual SAC-H-infinity gives the lowest final error.\label{fig:backbone_overall_sse}]{
\includegraphics[width=0.88\columnwidth]{overall_backbone_sse.png}}
\caption{Overall backbone comparison across the twenty-four benchmark cases.}
\label{fig:backbone_overall_pair}
\end{figure}

\begin{figure}[!b]
\centering
\subfloat[Hard-case backbone comparison: mean IAE. H-infinity retains the best accumulated-error behavior.\label{fig:backbone_hard_iae}]{
\includegraphics[width=0.88\columnwidth]{hard_backbone_iae.png}}\\[-1mm]
\subfloat[Hard-case backbone comparison: mean SSE. Residual SAC-H-infinity has the best final-error behavior.\label{fig:backbone_hard_sse}]{
\includegraphics[width=0.88\columnwidth]{hard_backbone_sse.png}}
\caption{Backbone comparison on the hard nonlinear subset.}
\label{fig:backbone_hard_pair}
\end{figure}
The case-by-case winner counts in Fig.~\ref{fig:backbone_winner_counts} further clarify the result. H-infinity wins IAE in nineteen of the twenty-four cases, while residual SAC-H-infinity wins IAE in five cases. For SSE, the result is split: H-infinity wins twelve cases and residual SAC-H-infinity wins twelve cases. Therefore, the contribution of residual SAC-H-infinity is not lower total error everywhere; it is improved final convergence on many cases while staying close to the best robust classical controller in accumulated tracking error.

\begin{figure}[!htbp]
\centering
\includegraphics[width=\columnwidth]{backbone_winner_counts.png}
\caption{Winner counts for the LQR, MPC, and H-infinity backbone study.\protect\\ H-infinity dominates IAE wins, while H-infinity and residual SAC-H-infinity split SSE wins.}
\label{fig:backbone_winner_counts}
\end{figure}

Fig.~\ref{fig:sac_mpc_training_curve} and Fig.~\ref{fig:sac_hinf_training_curve} show the training curves for residual SAC-MPC and residual SAC-H-infinity. As with residual SAC-LQR, these curves should be interpreted together with the curriculum. Reward does not increase monotonically because later stages contain harder nonlinear cases. The correct reading is that all three residual variants complete stable multi-seed training, after which the final controllers are selected using validation metrics rather than final raw reward. The SAC-MPC curves show that replacing LQR with MPC does not automatically make the residual policy stronger. The SAC-H-infinity curves show that a robust backbone gives the learned residual a better final-error operating point, but the robust classical controller itself remains very difficult to beat on IAE.

Taken together, this backbone study shows that the value of residual SAC is real but strongly dependent on the classical structure underneath it.

\begin{figure}[!t]
\centering
\subfloat[Residual SAC-MPC training curve. Replacing LQR with MPC does not automatically make the residual policy stronger.\label{fig:sac_mpc_training_curve}]{
\includegraphics[width=0.94\columnwidth]{sac_mpc_learning_curves.png}}\\[-1mm]
\subfloat[Residual SAC-H-infinity training curve. The robust backbone gives the learned residual a stronger final-error operating point.\label{fig:sac_hinf_training_curve}]{
\includegraphics[width=0.94\columnwidth]{sac_hinf_learning_curves.png}}
\caption{Training curves for residual SAC with alternative classical backbones.}
\label{fig:backbone_training_pair}
\end{figure}
\subsection{Discussion}
The results support three practical conclusions. First, pure model-free RL is not enough for this benchmark. Direct SAC and TD3 do not beat LQR or MPC on tracking, and direct SAC becomes overly conservative. Second, residual RL is much more competitive because it preserves the stabilizing behavior of a classical controller while allowing learned corrections. Third, the performance claim must be stated carefully. The proposed method is not ``pure SAC beats classical control.'' The correct conclusion is that residual SAC can improve the final-error behavior of strong classical controllers, but H-infinity state feedback remains the strongest controller on accumulated tracking error in this benchmark.

The energy result also deserves careful interpretation. Residual SAC-LQR has nearly the same energy use and saturation fraction as LQR. It is not a low-energy controller in the way direct SAC is. Instead, it uses actuator authority similarly to LQR but shapes the residual action to reduce final error and improve ramp tracking. Put differently, the final hybrid controller is not winning by being passive. It is winning by using almost the same amount of control effort more effectively.

Taken together, the graphs and tables reveal four distinct controller behaviors. Direct SAC is energy-conservative but inaccurate. LQR and MPC are accurate nominal or constraint-aware baselines. H-infinity is the strongest robust classical baseline for accumulated tracking error. Residual SAC-H-infinity gives the best steady-state error, while residual SAC-LQR and residual SAC-MPC show that a learned residual can improve the corresponding backbone but does not automatically dominate every classical design. That combined numerical pattern is the central empirical result of the paper.

\subsection{Limitations}
This study is simulation-based. The results depend on the selected motor model, nonlinear perturbation magnitudes, reward design, and validation protocol. The final IAE improvement of residual SAC over LQR and MPC is small, and residual SAC does not beat H-infinity on accumulated error, so hardware validation and statistical testing over additional plants are required before making deployment claims. The study also evaluates a fixed benchmark set rather than a continuous distribution of real operating conditions. Finally, the residual architecture relies on a good nominal backbone; if the model-based controller is poor, the benefit may be reduced.

These limitations keep the claims appropriately narrow and focused on simulation-benchmark evidence rather than universal deployment performance.

\section{Conclusion}
This paper presented an IEEE-style benchmark of classical and reinforcement learning controllers for DC motor speed control. Direct RL controllers did not outperform strong classical baselines: direct SAC reduced energy but suffered large tracking error, and direct TD3 remained weaker than LQR and MPC. A residual SAC framework was then evaluated using LQR, MPC, and H-infinity backbones with the same MPC warm-start data, curriculum training, and validation-based checkpoint selection.

The final results show that residual SAC-LQR improves standalone LQR and residual SAC-MPC improves standalone MPC, but the stronger H-infinity baseline changes the interpretation. Standalone H-infinity achieves the best mean IAE, $15.6040$, while residual SAC-H-infinity achieves the best mean SSE, $5.9890 \pm 0.0524$. The main conclusion is therefore not that RL universally outperforms classical control. The stronger and more defensible conclusion is that RL is most convincing here as a residual adaptive layer: it can improve final convergence and compensate nonlinear mismatch, but accumulated tracking-error performance still depends strongly on the quality of the classical backbone.

This point also clarifies what is different about the present paper relative to the earlier work. The earlier studies established the direct-controller benchmark and the initial residual SAC-LQR result, but they did not separate clearly whether the observed benefit came from residual learning itself or from the particular choice of LQR as the backbone. The present paper answers that question by extending the same residual SAC pipeline to MPC and H-infinity and by comparing all three under one fixed benchmark. That extension turns the paper from a single-controller case study into a backbone-comparison study, and the resulting core finding is more precise: the value of residual RL is real, but it is conditional on the strength and structure of the classical controller underneath it.

In this way, the conclusion closes on a more careful claim: residual RL can improve selected classical backbones, but it does not remove the need for strong model-based control.

\begin{thebibliography}{00}
\bibitem{kalman1960}
R. E. Kalman, ``Contributions to the theory of optimal control,'' \emph{Boletin de la Sociedad Matematica Mexicana}, vol. 5, no. 2, pp. 102--119, 1960.

\bibitem{anderson1971}
B. D. O. Anderson and J. B. Moore, \emph{Linear Optimal Control}. Englewood Cliffs, NJ, USA: Prentice-Hall, 1971.

\bibitem{astrom1995}
K. J. Astrom and T. Hagglund, \emph{PID Controllers: Theory, Design, and Tuning}, 2nd ed. Research Triangle Park, NC, USA: ISA, 1995.

\bibitem{mayne2000}
D. Q. Mayne, J. B. Rawlings, C. V. Rao, and P. O. M. Scokaert, ``Constrained model predictive control: Stability and optimality,'' \emph{Automatica}, vol. 36, no. 6, pp. 789--814, 2000.

\bibitem{doyle1989}
J. C. Doyle, K. Glover, P. P. Khargonekar, and B. A. Francis, ``State-space solutions to standard H2 and H-infinity control problems,'' \emph{IEEE Transactions on Automatic Control}, vol. 34, no. 8, pp. 831--847, 1989.

\bibitem{zhou1996}
K. Zhou, J. C. Doyle, and K. Glover, \emph{Robust and Optimal Control}. Upper Saddle River, NJ, USA: Prentice-Hall, 1996.

\bibitem{armstrong1994}
B. Armstrong-Helouvry, P. Dupont, and C. C. de Wit, ``A survey of models, analysis tools and compensation methods for the control of machines with friction,'' \emph{Automatica}, vol. 30, no. 7, pp. 1083--1138, 1994.

\bibitem{wang2023tec}
Y. Wang, S. Fang, J. Hu, and D. Huang, ``A novel active disturbance rejection control of PMSM based on deep reinforcement learning for more electric aircraft,'' \emph{IEEE Transactions on Energy Conversion}, vol. 38, no. 3, pp. 2586--2598, 2023.

\bibitem{wang2023tpe}
Y. Wang, S. Fang, and J. Hu, ``Active disturbance rejection control based on deep reinforcement learning of PMSM for more electric aircraft,'' \emph{IEEE Transactions on Power Electronics}, vol. 38, no. 1, pp. 406--416, 2023.

\bibitem{silver2014}
D. Silver, G. Lever, N. Heess, T. Degris, D. Wierstra, and M. Riedmiller, ``Deterministic policy gradient algorithms,'' in \emph{Proc. 31st Int. Conf. Machine Learning}, 2014, pp. 387--395.

\bibitem{lillicrap2015}
T. P. Lillicrap \emph{et al.}, ``Continuous control with deep reinforcement learning,'' \emph{arXiv preprint arXiv:1509.02971}, 2015.

\bibitem{schulman2015}
J. Schulman, S. Levine, P. Abbeel, M. Jordan, and P. Moritz, ``Trust region policy optimization,'' in \emph{Proc. 32nd Int. Conf. Machine Learning}, 2015, pp. 1889--1897.

\bibitem{schulman2017}
J. Schulman, F. Wolski, P. Dhariwal, A. Radford, and O. Klimov, ``Proximal policy optimization algorithms,'' \emph{arXiv preprint arXiv:1707.06347}, 2017.

\bibitem{fujimoto2018}
S. Fujimoto, H. van Hoof, and D. Meger, ``Addressing function approximation error in actor-critic methods,'' in \emph{Proc. 35th Int. Conf. Machine Learning}, 2018, pp. 1587--1596.

\bibitem{haarnoja2018}
T. Haarnoja, A. Zhou, P. Abbeel, and S. Levine, ``Soft actor-critic: Off-policy maximum entropy deep reinforcement learning with a stochastic actor,'' in \emph{Proc. 35th Int. Conf. Machine Learning}, 2018, pp. 1861--1870.

\bibitem{kormushev2013}
P. Kormushev, S. Calinon, and D. G. Caldwell, ``Reinforcement learning in robotics: Applications and real-world challenges,'' \emph{Robotics and Autonomous Systems}, vol. 62, no. 9, pp. 1230--1240, 2013.

\bibitem{francoislavet2018}
V. Francois-Lavet, P. Henderson, R. Islam, M. G. Bellemare, and J. Pineau, ``An introduction to deep reinforcement learning,'' \emph{Foundations and Trends in Machine Learning}, vol. 11, no. 3--4, pp. 219--354, 2018.

\bibitem{recht2019}
B. Recht, ``A tour of reinforcement learning: The view from continuous control,'' \emph{Annu. Rev. Control Robot. Auton. Syst.}, vol. 2, pp. 253--279, 2019.

\bibitem{motorreview2024}
L. Roveda, N. Iannacci, and D. Piga, ``Reinforcement learning for motor control: A comprehensive review,'' \emph{arXiv preprint arXiv:2412.17936}, 2024.

\bibitem{lewis2012}
F. L. Lewis, D. Vrabie, and K. G. Vamvoudakis, ``Reinforcement learning and feedback control: Using natural decision methods to design optimal adaptive controllers,'' \emph{IEEE Control Systems Magazine}, vol. 32, no. 6, pp. 76--105, 2012.

\bibitem{fazel2018}
M. Fazel, R. Ge, S. M. Kakade, and M. Mesbahi, ``Global convergence of policy gradient methods for the linear quadratic regulator,'' in \emph{Proc. 35th Int. Conf. Machine Learning}, 2018, pp. 1467--1476.

\bibitem{johannink2019}
T. Johannink, S. Bahl, A. Nair, J. Luo, A. Kumar, M. Loskyll, J. Ojea, E. Solowjow, and S. Levine, ``Residual reinforcement learning for robot control,'' in \emph{Proc. IEEE Int. Conf. Robotics and Automation}, 2019, pp. 6023--6029.

\bibitem{tobin2017}
J. Tobin \emph{et al.}, ``Domain randomization for transferring deep neural networks from simulation to the real world,'' in \emph{Proc. IEEE/RSJ Int. Conf. Intelligent Robots and Systems}, 2017, pp. 23--30.

\bibitem{peng2018}
X. B. Peng, M. Andrychowicz, W. Zaremba, and P. Abbeel, ``Sim-to-real transfer of robotic control with dynamics randomization,'' in \emph{Proc. IEEE Int. Conf. Robotics and Automation}, 2018, pp. 3803--3810.

\bibitem{rl_dcmotor_dqn}
A. Lerch, ``A reinforcement learning based controller for optimal speed control of a DC motor using deep Q-network algorithm,'' M.S. thesis, Politecnico di Milano, Milan, Italy, 2022.

\bibitem{adaptive_pi_dc_motor}
J. Peluffo-Ord\'{o}\~{n}ez, D. Rosero-Montalvo, and B. Baykant Alagoz, ``Adaptive PI controller based on a reinforcement learning algorithm for speed control of a DC motor,'' \emph{Biomimetics}, vol. 8, no. 5, art. 434, 2023.

\bibitem{tufenkci2023}
S. Tufenkci, B. Baykant Alagoz, G. Kavuran, C. Yeroglu, N. Herencsar, and S. Mahata, ``A theoretical demonstration for reinforcement learning of PI control dynamics for optimal speed control of DC motors by using twin delayed deep deterministic policy gradient algorithm,'' \emph{Symmetry}, vol. 15, no. 2, art. 360, 2023.

\bibitem{poudel2022}
B. Poudel, T. Watson, and W. Li, ``Learning to control DC motor for micromobility in real time with reinforcement learning,'' \emph{arXiv preprint arXiv:2108.00138}, 2021.

\bibitem{anugula2021}
R. Anugula and S. P. K. Karri, ``Deep reinforcement learning based adaptive controller of DC electric drive for reduced torque and current ripples,'' in \emph{Proc. IEEE Int. Conf. Technology, Research, and Innovation for Betterment of Society}, 2021, pp. 1--6.

\bibitem{qin2018}
S. Qin, F. Wang, B. Li, and Y. Chen, ``Improve PID controller through reinforcement learning,'' in \emph{Proc. IEEE CSAA Guidance, Navigation and Control Conf.}, 2018, pp. 1--6.

\bibitem{yin2023td3pmsm}
F. Yin, X. Yuan, Z. Ma, and X. Xu, ``Vector control of PMSM using TD3 reinforcement learning algorithm,'' \emph{Algorithms}, vol. 16, no. 9, art. 404, 2023.

\bibitem{hu2025td3pmsm}
Y. Hu, Z. Han, M. Li, and Y. Li, ``Speed optimization of permanent magnet synchronous motor based on deep reinforcement learning under twin delayed deep deterministic policy gradient,'' \emph{Energies}, vol. 18, no. 4, art. 901, 2025.

\bibitem{schindler2023iemdc}
A. Schindler, J. F. Staudinger, and M. J. M. Ruderman, ``Deep reinforcement learning current control of permanent magnet synchronous machines,'' in \emph{Proc. IEEE Int. Electric Machines \& Drives Conf.}, 2023, pp. 1--7.

\bibitem{zhao2023pmsm}
Z. Zhou, Y. Xu, and W. Wang, ``Reinforcement learning and optimal control of permanent magnet synchronous motor speed servo system,'' \emph{IEEE Transactions on Industrial Electronics}, vol. 70, no. 8, pp. 8074--8085, 2023.

\bibitem{wfhu2023multiscenario}
W.-F. Hu, M. Lv, J. Si, and C.-M. Lai, ``A multi-scenarios parameter optimization method based on deep reinforcement learning for permanent magnet synchronous motors,'' \emph{World Electric Vehicle Journal}, vol. 14, no. 8, art. 210, 2023.

\bibitem{dukkipati2026}
S. Dukkipati, J. Wang, Y. Wang, and K. Yuan, ``An adaptive deep reinforcement learning-based control strategy of permanent magnet synchronous motor for more electric aircraft,'' \emph{IEEE Transactions on Transportation Electrification}, vol. 12, no. 1, pp. 1--12, 2026.

\bibitem{review2025pmsm}
L. A. Moulai, F. M. Zaihidee, S. Mekhilef, J. R. Tang, and M. Mubin, ``Implementation of deep reinforcement learning in permanent magnet synchronous motors control: A review,'' \emph{Annual Reviews in Control}, vol. 60, art. 101014, 2025.

\end{thebibliography}

\end{document}
